\section{Workstream 4b: the two language models}
\label{sec:lm}

Both models use the existing \code{lm/model.py} (RoPE, RMSNorm, SwiGLU, PyTorch's fused attention)
and \code{lm/train.py} (bf16 autocast, AdamW, cosine schedule, checkpoint every 30 minutes, Ctrl+C
saves, resume on restart). Neither starts from anyone else's weights or tokenizer.

\subsection{TxGPT: a transaction-sequence model}
\label{sec:txgpt}

\textbf{Idea.} A language model learns which next ``word'' is likely. Here the words are the 6
tokens of each transaction (Section~\ref{sec:txtokens}), so the model learns how Bitcoin
addresses usually behave over time. A transaction the model finds unlikely, \emph{given that
address's own history}, is unusual. Because the model is causal, a transaction is scored only
from what came before it, which is time-safe by construction.

\textbf{Three sizes, to find the right one for the data.} DS2 gives about 6M training
tokens\est{}. That is small for a language model, so size is tuned rather than assumed:

\begin{center}
\small
\begin{tabularx}{\textwidth}{@{}L{1.5cm}C{1.3cm}C{1.3cm}C{1.3cm}C{1.8cm}C{1.9cm}Y@{}}
\toprule
\textbf{Size} & \textbf{Layers} & \textbf{$d$} & \textbf{Heads} & \textbf{SwiGLU hidden} & \textbf{Params\est{}} & \textbf{Peak learning rate} \\
\midrule
S & 4 & 128 & 4 & 352 & 0.8M & $1\times10^{-3}$ \\
M & 6 & 256 & 8 & 688 & 4.7M & $1\times10^{-3}$ \\
L & 8 & 384 & 6 & 1,024 & 14M & $6\times10^{-4}$ \\
\bottomrule
\end{tabularx}
\end{center}

\textbf{Training configuration} (\code{configs/txgpt_{s,m,l}.json}): sequence 768 tokens; batch
32 sequences (about 24,600 tokens); up to 20 passes over the data; validation every 100 steps;
early stopping after 5 validations without improvement; warm-up 100 steps; cosine decay to 10\%;
weight decay 0.1; dropout 0.1, because the data is small; bf16. Each size is trained with 3 seeds,
9 runs in total, at minutes each\est{}. The size with the lowest mean validation loss wins. A
plot of validation loss against parameter count goes in the appendix as evidence of tuning.

\textbf{Is it learning anything? Perplexity against simple baselines.} On the held-out test
addresses, compare the average negative log-likelihood per token and per transaction, field by
field, with three baselines that need no neural network:
\begin{enumerate}
  \item \textbf{Global frequencies}: each field's overall token frequency.
  \item \textbf{Address habit}: each field's token frequency in \emph{this address's} history so
    far (Laplace-smoothed). This is close to what the robust-z baseline assumes.
  \item \textbf{First-order Markov}: each field's token given the same field in the previous
    transaction.
\end{enumerate}
If TxGPT beats the address-habit and Markov baselines, it has learned structure beyond counting.
If it doesn't, that is reported plainly as a finding, and TxGPT is not chosen.

\textbf{Scoring DS1} (\code{python -m lm.tx_score}):
\begin{enumerate}
  \item Tokenise each DS1 address's full history with DS2's bucket edges (from \code{meta.json}).
  \item Slide 768-token windows so that every transaction is scored with as much earlier context
    as fits (up to 127 earlier transactions).
  \item A transaction's surprisal is $s_i = -\sum_{f\in\text{fields}} \log p(\text{token}_f \mid
    \text{everything before it})$, kept per field.
  \item Two candidate scores: raw $s_i$, and $s_i$ relative to the address (a robust z against
    the address's own earlier surprisals). The tuning plant set chooses between them.
  \item Apply the same minimum-history rule as the baseline (at least 10 earlier transactions
    over at least 14 days), so the methods are compared on the same rows. The flag threshold is
    the top 2\% of fit-period scores.
  \item \textbf{Explanation.} The field with the largest surprisal names the reason: gap →
    ``the timing was unusual for this address'', amount → ``the amount was unusual'', inputs and
    outputs → ``the transaction's shape was unusual'', fee → ``the fee rate was unusual''.
\end{enumerate}

\textbf{Export.} \code{models/txgpt.pt} (fp16 weights, about 30 MB for size L\est{}) plus its
\code{meta.json}. Scoring runs on CPU in the marker's environment, within minutes for DS1.

\subsection{The Explorer Assistant: a 152M text model}
\label{sec:assistant}

\textbf{Stage 1, pretraining.} The existing \code{configs/text_150m_pretrain.json}: 18 layers,
$d=768$, 12 heads, SwiGLU hidden 2,048, sequence 1,024, our 32,768-token BPE tokenizer. 2B tokens
of DS5, subject to the gate in Section~\ref{sec:gate}. Batch about 262,000 tokens; peak learning
rate $6\times10^{-4}$ falling to $6\times10^{-5}$; warm-up 500 steps; weight decay 0.1. It starts
on the evening of 21 October.

\textbf{Stage 2, chat fine-tune.} The existing \code{configs/text_150m_chat.json}: 100M tokens of
DS6. Loss only on assistant replies (the mask file); learning rate $2\times10^{-4}$ falling to
$2\times10^{-5}$; about 3--4 hours\est{}.

\textbf{Stage 3, explanation fine-tune (new).} \code{configs/text_150m_explain.json}: start from
the chat model; DS7 for 3 passes (about 8M tokens\est{}); learning rate $1\times10^{-4}$; warm-up
20 steps; loss on replies only; validation on DS7's validation split. Under an hour\est{}.

\textbf{Evaluation.}
\begin{center}
\small
\begin{tabularx}{\textwidth}{@{}L{4.0cm}Y@{}}
\toprule
\textbf{Measure} & \textbf{How} \\
\midrule
Language-modelling quality & Validation loss and perplexity on DS5's 10M held-out tokens after stage 1, on DS6 validation after stage 2, and on DS7 validation after stage 3. Training and validation curves from each run's \code{log.jsonl}. \\
Number faithfulness & On about 500 DS7 test inputs with greedy decoding: the share of outputs in which \emph{every} number appears in the input, after formatting. \\
Limitation line & The share of outputs containing the fixed line exactly. \\
Banned words & The share of outputs containing any of: fraud, criminal, illicit, guilty, suspicious. The target is 0 after the validator. \\
Feature coverage & The share of contributing features the explanation mentions. \\
Similarity to the template & ROUGE-L against the template answer, written as a short LCS function rather than another dependency. \\
Human check & A teammate rates 30 random outputs for clarity and correctness (1--5), blind to which model wrote them; 10 template outputs are mixed in as controls. \\
\bottomrule
\end{tabularx}
\end{center}

\textbf{Guardrail for A3.} \code{explorer/assistant.py} runs the model, then validates the output:
the numbers must match the input, the limitation line must be present, and no banned word may
appear. If any check fails, it returns the template explanation instead. The pass rate is
reported. This keeps the Explorer trustworthy even if the small model hallucinates, and the
report says so: ``the template is authoritative; the assistant only rephrases''.

\textbf{Honest framing for the report.} These are \emph{small} language models trained from
scratch on one consumer GPU. They are not comparable to commercial LLMs. Their value is that
they show the whole pipeline: tokenizer, pretraining, alignment, evaluation and a safety wrapper.

\subsection{The Explorer Assistant 3B: a fine-tuned comparison (added 22 October)}
\label{sec:assistant3b}

The tutor, Arthur Hoang, agreed to an extra model: a pretrained base model of about 3B parameters
(Qwen2.5-3B-Instruct, revision \code{aa8e7253}, chosen on 23 October: not gated, and already chat-tuned, so it can also serve as the Explorer's chatbot), fine-tuned by the team. It is added
to the 152M Assistant and doesn't replace it.

\textbf{Method.} QLoRA: the base weights are loaded in 4-bit and only low-rank adapters are
trained, which fits the RTX 3060's 12\,GB. One stage: DS7 mixed with 2,500 DS6 conversations
(so its general chat does not degrade), two passes, loss on replies only; about 10.5M tokens. The DS7 train, validation and
test splits are the same ones the 152M model uses. A few hours of GPU time in total\est{}.

\textbf{When.} After the 152M model's three stages and TxGPT are done (both have priority).
It never pauses those runs.

\textbf{Evaluation and guardrail.} The same measures as the table above, on the same DS7 test
inputs, so the report has a direct from-scratch vs fine-tuned comparison. It serves 4-bit
quantised behind the same \code{explorer/assistant.py} validator, with the template answer as the
fallback.

\textbf{Records.} The base model's name, revision and licence go in \code{LICENSES.md} and the
model card. The change goes in \code{CHANGELOG.md} and the AI declaration. The adapter weights
ship (or a link to them if they are too large); the base weights don't.

\subsection{GPU timetable}

\begin{center}
\small
\begin{tabularx}{\textwidth}{@{}L{3.4cm}Y@{}}
\toprule
\textbf{When} & \textbf{GPU job} \\
\midrule
Wed 21 Oct, evening & Benchmark → apply the gate → start pretraining. \\
Thu 22 -- Sun 25 Oct & Pretraining. Check loss and temperature twice a day. \\
Mon 26 Oct & Pretraining ends (or is paused if it overruns). TxGPT sweep, 9 runs. \\
Tue 27 Oct & TxGPT scoring of DS1; perplexity baselines. Resume pretraining if it was paused. \\
Wed 28 Oct & Chat fine-tune, 3--4 h. \\
Thu 29 -- Fri 30 Oct & Build DS7 → explanation fine-tune → faithfulness evaluation. TxGPT plant evaluation (CPU or GPU). \\
After 30 Oct & 3B QLoRA fine-tune (chat, then DS7) and its evaluation, a few hours. Then GPU free; re-runs only if a bug is found. \\
\bottomrule
\end{tabularx}
\end{center}

\begin{caution}
If pretraining hasn't finished by Monday 26 October even after the gate, stop it there. Run
stages 2 and 3 from the latest checkpoint and report the token count actually reached. TxGPT is
the HD-critical model and must not wait for the text model.
\end{caution}

\begin{feeds}
Report §6: two short paragraphs justifying TxGPT (sequence context, time-safe, explainable) and
the Assistant (the tutor's requirement; an explanation layer behind a guardrail). Report §7.2:
the perplexity table against the baselines, TxGPT's plant recall in the main comparison, and the
Assistant's faithfulness table. Appendix: configs, training curves, the scaling plot and sample
outputs. HD ``additional ML'': both models.
\end{feeds}
